\documentclass[11pt,a4paper]{article}
\usepackage[T1]{fontenc}
\usepackage{lmodern,amsmath,amssymb,amsthm,mathtools,microtype}
\usepackage[margin=27mm]{geometry}
\usepackage[colorlinks=true,linkcolor=blue,citecolor=blue,urlcolor=blue]{hyperref}
\setlength{\emergencystretch}{2em}
\newtheorem{theorem}{Theorem}[section]
\newtheorem{corollary}[theorem]{Corollary}
\newtheorem{remark}[theorem]{Remark}
\newcommand{\Q}{\mathbb Q}
\newcommand{\N}{\operatorname N}
\newcommand{\Tr}{\operatorname{Tr}}
\title{Doubled source points on rational congruent-number twists}
\author{}
\date{16 November 2026}
\begin{document}
\maketitle
\begin{abstract}
A nonzero doubled point on $y^2=x^3-(3-\sqrt2)^2x$ gives an explicit
rational point on a congruent-number twist. Its abscissa has the square
class of the weighted trace $3x_0+2x_1$, where
$x=x_0+x_1\sqrt2$. The construction uses only rational norms, and all
nonzero conditions follow from doubling and the two real embeddings.
\end{abstract}

\section{The norm construction}

Put $K=\Q(\sqrt2)$, $d=3-\sqrt2$, and
\[
                         E:\quad y^2=x^3-d^2x.
\]
A bar denotes the nontrivial automorphism of $K/\Q$.

\begin{theorem}
Let $R=(x,y)\in2E(K)\setminus\{O\}$. Choose $h,j\in K$ with
$x=h^2$ and $j^2=x^2-d^2$, and put
\[
 a=\N(h),\qquad b=\N(j),\qquad
 r=\tfrac12\Tr(x\bar d),\qquad L=\N(x+d),\qquad m=rb.
\]
Then $r>0$, $b\ne0$, $L$ is a nonzero rational square, and
\[
 b^2=(a^2+7)^2-4r^2,\qquad L=a^2+7+2r.
\]
The rational point
\[
                         (X,Y)=(rL,2r^2L)
\]
has nonzero ordinate on the congruent-number curve
\[
                         Y^2=X^3-m^2X,
\]
and its abscissa has rational square class $[X]=[r]$.
\end{theorem}
\begin{proof}
Write $R=2Q$ with $Q=(z,w)$. Since $R\ne O$, the point $Q$ is finite
and $w\ne0$. The duplication identities give
\begin{align*}
 h&=\frac{z^2+d^2}{2w},&
 x&=h^2,\\
 k_+&=\frac{z^2+2dz-d^2}{2w},&
 x+d&=k_+^2,\\
 k_-&=\frac{z^2-2dz-d^2}{2w},&
 x-d&=k_-^2.
\end{align*}
Thus $j=k_+k_-$ is an available choice. At either real embedding,
$z^2+d^2>0$, so $h\ne0$ and $x>0$. Moreover $x\ne\pm d$: otherwise
$\N(x)=\N(h)^2=7$, an impossibility for rational $\N(h)$.
Hence $j\ne0$, and $x^2-d^2=j^2$ implies $x>d>0$ at both embeddings.
In particular $r>0$, $b\ne0$, and $L=\N(k_+)^2>0$.

Since $\N(d)=7$ and $\N(x)=a^2$, direct multiplication gives
\begin{align*}
 \N(x+d)&=a^2+7+x\bar d+\bar x d=a^2+7+2r,\\
 \N(x-d)&=a^2+7-2r.
\end{align*}
Taking norms of $j^2=(x-d)(x+d)$ proves
\[
             b^2=(a^2+7-2r)(a^2+7+2r).
\]
It follows that $L^2-b^2=4rL$, and therefore
\[
 (rL)^3-(rb)^2(rL)=r^3L(L^2-b^2)=4r^4L^2=(2r^2L)^2.
\]
The nonzero and square-class assertions follow from $r,L>0$ and
$L\in\Q^{\times2}$.
\end{proof}

\begin{corollary}
Every such point produces a rational right triangle with nonzero legs
$|b|$ and $2r$, hypotenuse $a^2+7$, and area $|m|$.
\end{corollary}

\section{An exact source point}

For the doubled source point
\[
 P=2(1+\sqrt2,2\sqrt2)
   =\left(\frac{57-28\sqrt2}{8},
            \frac{-660+417\sqrt2}{32}\right),
\]
one obtains
\[
  a=-\frac{41}{8},\quad b=-\frac{1071}{64},\quad
  r=\frac{115}{8},\quad L=\frac{3969}{64},\quad
  m=-\frac{123165}{512}.
\]
The triangle identity is
$1071^2+1840^2=2129^2$ after clearing the denominator $64$.

\begin{remark}
The construction prescribes the square class of the weighted trace
$r=3x_0+2x_1$ when $x=x_0+x_1\sqrt2$. It does not identify this class
with an independently prescribed rational input class, nor does it impose
a complementary twist relation.
\end{remark}

The script
\nolinkurl{computation/pole_parameters/doubled_source_congruent_twists_v43/verify.py}
verifies the norm and duplication identities symbolically and checks
this source point with exact arithmetic in $\Q(\sqrt2)$.
\end{document}
